import matplotlib.pyplot as plt
from sklearn.cluster import KMeans

# WCSS-Werte für verschiedene k berechnen
wcss = []
k_range = range(1, 11)

for k in k_range:
    kmeans = KMeans(n_clusters=k, random_state=42, n_init=10)
    kmeans.fit(wine.data)
    wcss.append(kmeans.inertia_)

# Elbow-Plot erstellen
plt.figure(figsize=(8, 6))
plt.plot(k_range, wcss, 'o-', color='blue', linewidth=2, markersize=8)
plt.xlabel('Anzahl Cluster (k)', fontsize=12)
plt.ylabel('WCSS (Within-Cluster Sum of Squares)', fontsize=12)
plt.title('Elbow-Methode für den Wine-Datensatz', fontsize=14)
plt.xticks(k_range)
plt.grid(True, alpha=0.3)

# Knick bei k=3 markieren
plt.axvline(x=3, color='red', linestyle='--', alpha=0.7, label='Elbow bei k=3')
plt.legend()

plt.tight_layout()
plt.show()